\chapter{二次函数与二次方程}\label{ch:g11:quad}

\omterm{def:g11:quad:quadratic}{二次函数}是继\omterm{def:g10:reffunc:affine}{一次函数}之后最简单的\omterm{def:g10:functions:function}{函数}，也是
首个\omterm{def:g10:functions:graph}{图像}真正弯曲的\omterm{def:g10:functions:function}{函数}。本章建立关于它们的完整工具：
\omterm{thm:g11:quad:canonical}{顶点式}、\omterm{def:g11:quad:discriminant}{判别式}、二次式的符号，以及
\omterm{def:g11:quad:parabola}{抛物线}的几何性质。

\section{顶点式}

\begin{definition}[二次函数]\label{def:g11:quad:quadratic}
形如下式的\omterm{def:g10:functions:function}{函数}称为\emph{二次函数}\index{二次函数}，
即
\[
f(x) = ax^2 + bx + c, \qquad a \neq 0,
\]
其中 $a$、$b$、$c$ 为\omterm{def:g10:numbers:sets}{实数}。式子 $ax^2 + bx + c$
也称为\emph{二次三项式}。
\end{definition}

\begin{theorem}[顶点式]\label{thm:g11:quad:canonical}
每个\omterm{def:g11:quad:quadratic}{二次函数}都能唯一地写成
\[
f(x) = a(x - \alpha)^2 + \beta,
\qquad\text{其中 } \alpha = -\frac{b}{2a} \text{ 且 } \beta = f(\alpha).
\]
\emph{顶点式}\index{顶点式}即 $f$ 的上述表达式。
\end{theorem}

\begin{proof}
用\emph{配方法}。由 $a \neq 0$，
将其提出前两项：
\[
ax^2 + bx + c = a\left(x^2 + \frac{b}{a}x\right) + c .
\]
现在 $x^2 + \frac{b}{a}x$ 是平方式的前两项，
$\left(x + \frac{b}{2a}\right)^2 = x^2 + \frac ba x + \frac{b^2}{4a^2}$,
故 $x^2 + \frac{b}{a}x = \left(x + \frac{b}{2a}\right)^2 -
\frac{b^2}{4a^2}$。代入得：
\[
f(x) = a\left(x + \frac{b}{2a}\right)^2 - \frac{b^2}{4a} + c
= a\bigl(x - \alpha\bigr)^2 + \beta,
\]
其中 $\alpha = -\frac{b}{2a}$，$\beta = c - \frac{b^2}{4a}$。
取 $x = \alpha$，平方项为零，故 $\beta = f(\alpha)$。唯一性：
将 $a(x-\alpha)^2 + \beta$ \omterm{def:g10:algebra:expand}{展开}，比较
$x$ 的系数可得 $\alpha = -\frac{b}{2a}$，进而 $\beta$ 被唯一确定。
\end{proof}

\begin{example}\label{ex:g11:quad:canonical}
设 $f(x) = 2x^2 - 8x + 3$，则 $\alpha = -\frac{-8}{2 \times 2} = 2$，且
$\beta = f(2) = 8 - 16 + 3 = -5$，故
\[
f(x) = 2(x - 2)^2 - 5 .
\]
\omterm{def:g10:algebra:expand}{展开}验证：$2(x^2 - 4x + 4) - 5 = 2x^2 - 8x + 8 - 5 =
2x^2 - 8x + 3$.
\end{example}

\begin{proposition}[极值与单调性]\label{prop:g11:quad:variations}
设 $f(x) = a(x-\alpha)^2 + \beta$。
\begin{enumerate}
  \item 若 $a > 0$，则 $f$ 在 $\intoc{-\infty}{\alpha}$ 上\omterm{def:g10:functions:variations}{递减}，且
        在 $\intco{\alpha}{+\infty}$ 上\omterm{def:g10:functions:variations}{递增}：$f$ 有\emph{\omterm{def:g10:functions:extrema}{最小值}}
        $\beta$，在 $x = \alpha$ 处取得。
  \item 若 $a < 0$，则单调性相反：$f$ 有\emph{\omterm{def:g10:functions:extrema}{最大值}}
        $\beta$，在 $x = \alpha$ 处取得。
\end{enumerate}
\end{proposition}

\begin{proof}
设 $a > 0$，取 $\alpha \leq u < v$。则
\[
f(v) - f(u) = a\bigl((v-\alpha)^2 - (u-\alpha)^2\bigr)
= a\,(v - u)\,\bigl((v-\alpha) + (u-\alpha)\bigr).
\]
各因式均为正：$a > 0$，$v - u > 0$，且
$(v-\alpha) + (u-\alpha) > 0$（因为 $v > \alpha$，$u \geq \alpha$）。于是
$f(v) > f(u)$，即 $f$ 在 $\intco{\alpha}{+\infty}$ 上\omterm{def:g10:functions:variations}{递增}。对于
$u < v \leq \alpha$ 时，最后一个因式为负，故 $f(v) < f(u)$，即 $f$
在该\omterm{def:g10:numbers:interval}{区间}上\omterm{def:g10:functions:variations}{递减}。最后，$f(x) - f(\alpha) = a(x-\alpha)^2 \geq 0$ 对
所有 $x$ 成立，故 $\beta = f(\alpha)$ 为\omterm{def:g10:functions:extrema}{最小值}。$a < 0$ 的情形由
上述结论应用于 $-f$ 即可得出。
\end{proof}

\begin{example}
某农夫有 $100$ 米篱笆，用来围一块紧靠直墙的矩形田地
（靠墙一侧无需篱笆）。若 $x$ 为宽度，则
围成面积为 $A(x) = x(100 - 2x) = -2x^2 + 100x$。此时 $\alpha =
-\frac{100}{2\times(-2)} = 25$，且 $A(25) = 25 \times 50 = 1250$：面积
在宽为 $25$ m 时最大，为 $1250$ m$^2$。
\end{example}

\section{根与因式分解}

\begin{definition}[根、判别式]\label{def:g11:quad:discriminant}
\emph{根}\index{根}指\omterm{def:g11:quad:quadratic}{二次三项式} $ax^2+bx+c$ 所对应的
\omterm{def:g10:algebra:equation}{方程} $ax^2 + bx + c = 0$ 的解。\omterm{def:g11:quad:quadratic}{二次三项式}的\emph{判别式}\index{判别式}
为如下的数：
\[
\Delta = b^2 - 4ac .
\]
\end{definition}

\begin{theorem}[解方程 $ax^2+bx+c=0$]\label{thm:g11:quad:roots}
\begin{enumerate}
  \item 若 $\Delta > 0$，则\omterm{def:g10:algebra:equation}{方程}有两个不相等的根
        \[
        x_1 = \frac{-b - \sqrt{\Delta}}{2a},
        \qquad
        x_2 = \frac{-b + \sqrt{\Delta}}{2a},
        \]
        且 $ax^2+bx+c = a(x - x_1)(x - x_2)$。
  \item 若 $\Delta = 0$，则\omterm{def:g10:algebra:equation}{方程}有唯一的根
        $x_0 = -\frac{b}{2a}$，且 $ax^2+bx+c = a(x - x_0)^2$。
  \item 若 $\Delta < 0$，则\omterm{def:g10:algebra:equation}{方程}无\omterm{def:g10:numbers:sets}{实数}解，且
        该\omterm{def:g11:quad:quadratic}{二次三项式}不能分解为实系数的因式。
\end{enumerate}
\end{theorem}

\begin{proof}
由\cref{thm:g11:quad:canonical}中算出的\omterm{thm:g11:quad:canonical}{顶点式}出发：
\[
ax^2+bx+c
= a\left(x + \frac{b}{2a}\right)^2 + c - \frac{b^2}{4a}
= a\left[\left(x + \frac{b}{2a}\right)^2 - \frac{\Delta}{4a^2}\right],
\]
因为 $c - \frac{b^2}{4a} = \frac{4ac - b^2}{4a} = -\frac{\Delta}{4a} =
a \times \left(-\frac{\Delta}{4a^2}\right)$。

\emph{1.} 若 $\Delta > 0$，则 $\frac{\Delta}{4a^2} =
\left(\frac{\sqrt\Delta}{2a}\right)^2$，且括号部分是
两个平方之差：
\[
ax^2+bx+c = a\left(x + \frac{b}{2a} - \frac{\sqrt\Delta}{2a}\right)
\left(x + \frac{b}{2a} + \frac{\sqrt\Delta}{2a}\right)
= a(x - x_2)(x - x_1).
\]
乘积为零，当且仅当其中某个因式为零，由此得出
两个根，因 $\sqrt\Delta \neq 0$，二者互异。

\emph{2.} 若 $\Delta = 0$，则括号部分为
$\left(x + \frac{b}{2a}\right)^2$，它仅在
$x_0 = -\frac{b}{2a}$ 处为零。

\emph{3.} 若 $\Delta < 0$，则 $-\frac{\Delta}{4a^2} > 0$，因此括号部分
是一个平方加一个正数，故永不为零，\omterm{def:g10:algebra:equation}{方程}
无\omterm{def:g10:numbers:sets}{实数}解。若能分解为 $a(x-x_1)(x-x_2)$，则会产生根，因此
这样的分解不存在。
\end{proof}

\begin{example}\label{ex:g11:quad:solving}
解\omterm{def:g10:algebra:equation}{方程} $2x^2 - 8x + 3 = 0$。此时 $\Delta = (-8)^2 - 4 \times 2 \times 3 =
64 - 24 = 40 > 0$，故\omterm{def:g10:algebra:equation}{方程}有两个根：
\[
x_{1,2} = \frac{8 \pm \sqrt{40}}{4} = \frac{8 \pm 2\sqrt{10}}{4}
= 2 \pm \frac{\sqrt{10}}{2}.
\]
对于 $x^2 + x + 1 = 0$，$\Delta = 1 - 4 = -3 < 0$，无\omterm{def:g10:numbers:sets}{实数}解。
\end{example}

\begin{proposition}[根的和与积]\label{prop:g11:quad:vieta}
若 $\Delta \geq 0$，且 $x_1, x_2$ 为其根（$\Delta = 0$ 时二者相等），
则
\[
x_1 + x_2 = -\frac{b}{a},
\qquad
x_1 x_2 = \frac{c}{a} .
\]
反之，和为 $s$、积为 $p$ 的两个数，即
$x^2 - sx + p = 0$ 的两个解。
\end{proposition}

\begin{proof}
将\cref{thm:g11:quad:roots}中的\omterm{def:g10:algebra:expand}{因式分解}\omterm{def:g10:algebra:expand}{展开}：
\[
a(x - x_1)(x - x_2) = a x^2 - a(x_1 + x_2)x + a\,x_1 x_2 .
\]
比较 $ax^2 + bx + c$ 的系数，得
$-a(x_1+x_2) = b$ 与 $a\,x_1x_2 = c$。反之，若 $u + v = s$ 且
$uv = p$，则 $(x-u)(x-v) = x^2 - sx + p$，故 $u$、$v$ 为其根。
\end{proof}

\begin{example}
求和为 $10$、积为 $21$ 的两个数。它们满足
$x^2 - 10x + 21 = 0$；$\Delta = 100 - 84 = 16$，根为
$\frac{10 \pm 4}{2} = 7$ 与 $3$。\emph{快速找根：}由于
$x^2 - 5x + 4$ 的系数满足 $1 - 5 + 4 = 0$，故数 $1$ 是一个
根，另一根为 $\frac{c}{a} = 4$（两根之积必为
$\frac ca$）。
\end{example}

\section{二次三项式的符号}

\begin{theorem}[二次三项式的符号]\label{thm:g11:quad:sign}
设 $f(x) = ax^2+bx+c$，其中 $a \neq 0$。
\begin{enumerate}
  \item 若 $\Delta > 0$，根为 $x_1 < x_2$，则 $f(x)$ 与 $a$ 同号的
        部分是 $\intcc{x_1}{x_2}$ 之外，与 $-a$ 同号的
        部分是 $\intoo{x_1}{x_2}$。
  \item 若 $\Delta = 0$，则 $f(x)$ 与 $a$ 同号，对一切 $x \neq x_0$
        成立，且在 $x_0$ 处为零。
  \item 若 $\Delta < 0$，则 $f(x)$ 与 $a$ 同号，对一切 $x$ 成立。
\end{enumerate}
简言之：\emph{\omterm{def:g11:quad:quadratic}{二次三项式}与 $a$ 同号，两根之间除外}。
\end{theorem}

\begin{proof}
\emph{1.} 写 $f(x) = a(x-x_1)(x-x_2)$，并考察各因式的
符号。若 $x < x_1$，则 $x - x_1$ 与 $x - x_2$ 均为负，其
积为正，故 $f(x)$ 与 $a$ 同号。若 $x_1 < x < x_2$，两因式
异号，积为负，故 $f(x)$ 与 $-a$ 同号。若 $x > x_2$，两因式
均为正，故 $f(x)$ 仍与
$a$ 同号。

\emph{2.} $f(x) = a(x - x_0)^2$，且 $(x-x_0)^2 > 0$，只要 $x \neq x_0$。

\emph{3.} 由\cref{thm:g11:quad:roots}的证明，
$f(x) = a\left[\left(x + \frac{b}{2a}\right)^2 +
\left(-\frac{\Delta}{4a^2}\right)\right]$，括号部分
恒为正。
\end{proof}

\begin{method}[解二次不等式]\label{met:g11:quad:inequalities}
求解 $ax^2 + bx + c \leq 0$（或 $\geq$、$<$、$>$），步骤如下：
\begin{enumerate}
  \item 计算 $\Delta$，若有根则求出根；
  \item 利用\cref{thm:g11:quad:sign}列出\omterm{def:g10:algebra:signtable}{符号表}（
        根之外取 $a$ 的符号）；
  \item 写出解\omterm{def:g10:numbers:interval}{区间}，并注意根本身是否属于解集
        （非严格不等式时属于，严格不等式时不属于）。
\end{enumerate}
\end{method}

\begin{example}\label{ex:g11:quad:inequality}
求解 $-x^2 + 3x + 4 \geq 0$。这里 $\Delta = 9 + 16 = 25$，根为
$\frac{-3 \pm 5}{-2}$，即 $x_1 = -1$ 与 $x_2 = 4$。由于
$a = -1 < 0$，三项式在两根之外为负，
在两根之间为正。由于是非严格不等式，解集为
$\intcc{-1}{4}$。
\end{example}

\section{抛物线}

\begin{definition}[抛物线]\label{def:g11:quad:parabola}
\omterm{def:g11:quad:quadratic}{二次函数} $f(x) = a(x-\alpha)^2+\beta$ 的\omterm{def:g10:functions:graph}{图像}是一条
\emph{抛物线}\index{抛物线}，其\emph{顶点}为 $S(\alpha, \beta)$。
当 $a > 0$ 时开口向上，$a < 0$ 时开口向下，且关于
竖直直线 $x = \alpha$ 对称。
\end{definition}

\begin{remark}
对称性易于验证：对任意 $h$，
$f(\alpha + h) = ah^2 + \beta = f(\alpha - h)$，因此与直线 $x = \alpha$ 水平距离相等的两点
高度相同。
\end{remark}

\begin{omfigure}
\begin{tikzpicture}
\begin{axis}[omaxis,
  width=5.1cm, height=4.8cm,
  xmin=-0.8, xmax=4.4, ymin=-1.7, ymax=3.6,
  xtick={1,3}, xticklabels={$x_1$,$x_2$}, ytick=\empty,
  title={$\Delta > 0$}]
  \addplot[omDef, thick, domain=-0.15:4.15] {(x-1)*(x-3)};
  \fill (axis cs:1,0) circle (1.5pt);
  \fill (axis cs:3,0) circle (1.5pt);
  \fill (axis cs:2,-1) circle (1.5pt) node[below, font=\small] {$S$};
\end{axis}
\end{tikzpicture}%
\begin{tikzpicture}
\begin{axis}[omaxis,
  width=5.1cm, height=4.8cm,
  xmin=-0.8, xmax=4.4, ymin=-1.7, ymax=3.6,
  xtick={2}, xticklabels={$x_0$}, ytick=\empty,
  title={$\Delta = 0$}]
  \addplot[omDef, thick, domain=0.2:3.8] {(x-2)^2};
  \fill (axis cs:2,0) circle (1.5pt)
    node[above right=1pt, font=\small] {$S$};
\end{axis}
\end{tikzpicture}%
\begin{tikzpicture}
\begin{axis}[omaxis,
  width=5.1cm, height=4.8cm,
  xmin=-0.8, xmax=4.4, ymin=-1.7, ymax=3.6,
  xtick=\empty, ytick=\empty,
  title={$\Delta < 0$}]
  \addplot[omDef, thick, domain=0.35:3.65] {(x-2)^2 + 0.8};
  \fill (axis cs:2,0.8) circle (1.5pt) node[below, font=\small] {$S$};
\end{axis}
\end{tikzpicture}

{\small 开口向上的\omterm{def:g11:quad:parabola}{抛物线}（$a > 0$）相对于
$x$ 轴的三种位置：两个根、一个重根、无根。顶点 $S$ 的\omterm{def:g10:coordgeom:system}{横坐标}为
$\alpha = -\frac{b}{2a}$。}
\end{omfigure}

\begin{omfigure}
\begin{tikzpicture}
\begin{axis}[omaxis,
  width=8.6cm, height=5.6cm,
  xmin=-2.2, xmax=5.4, ymin=-4.4, ymax=6.2,
  xtick={-1,4}, xticklabels={$x_1$,$x_2$}, ytick=\empty]
  \addplot[name path=par, omDef, thick, domain=-1.75:4.75, samples=80]
    {-0.8*(x+1)*(x-4)};
  \path[name path=xaxis] (-1.75,0) -- (4.75,0);
  \addplot[omDef!20] fill between[of=par and xaxis,
    soft clip={domain=-1:4}];
  \fill (axis cs:-1,0) circle (1.5pt);
  \fill (axis cs:4,0) circle (1.5pt);
  \node[font=\small, omDef] at (axis cs:1.5,2.0) {$f > 0$};
  \node[font=\small, omThm] at (axis cs:-1.75,-1.3) {$f < 0$};
  \node[font=\small, omThm] at (axis cs:4.75,-1.3) {$f < 0$};
\end{axis}
\end{tikzpicture}

{\small 有两个根的开口向下三项式（$a < 0$）的符号：
两根之间（阴影部分）取 $-a$ 的符号，两根之外取 $a$ 的符号。}
\end{omfigure}

\begin{method}[抛物线的作图]\label{met:g11:quad:reading}
作 $f(x) = ax^2+bx+c$ 的\omterm{def:g10:functions:graph}{图像}：先确定开口方向（$a$ 的符号）、
顶点 $\left(-\frac{b}{2a},\, f\!\left(-\frac{b}{2a}\right)\right)$、
点 $(0, c)$ 为与 $y$ 轴的交点，以及根（若有）。
对称轴 $x = -\frac{b}{2a}$ 使每个已算出的点都在\omterm{def:g10:functions:graph}{图像}上出现两次。
\end{method}

\begin{example}
对 $f(x) = x^2 - 4x + 3$：开口向上，顶点 $(2, -1)$，与 $y$ 轴
相交于 $(0,3)$，根为 $1$ 与 $3$（观察可得 $1 - 4 + 3 = 0$）。由对称性，
点 $(4, 3)$ 也在\omterm{def:g10:functions:graph}{图像}上，是 $(0, 3)$ 的对称点。
\end{example}

\section{练习}

\begin{exercise}[$\star$]\label{exo:g11:quad:1}
解下列\omterm{def:g10:algebra:equation}{方程}
\[
x^2 - 5x + 6 = 0, \qquad
2x^2 + 3x - 2 = 0, \qquad
x^2 + 2x + 5 = 0, \qquad
9x^2 - 6x + 1 = 0 .
\]
\end{exercise}

\begin{exercise}[$\star$]\label{exo:g11:quad:2}
将下列\omterm{def:g10:functions:function}{函数}化为\omterm{thm:g11:quad:canonical}{顶点式}，并求\omterm{def:g11:quad:parabola}{抛物线}的顶点：
\[
f(x) = x^2 - 6x + 11, \qquad
g(x) = -2x^2 + 4x + 1, \qquad
h(x) = 3x^2 + 6x .
\]
\end{exercise}

\begin{exercise}[$\star$]\label{exo:g11:quad:3}
解下列不等式
\[
x^2 - 3x - 10 < 0, \qquad
2x^2 + x + 3 > 0, \qquad
-x^2 + 6x - 9 \geq 0 .
\]
\end{exercise}

\begin{exercise}[$\star$]\label{exo:g11:quad:4}
不计算 $\Delta$，先求出一个显然的根，再求另一个根：
\[
x^2 - 7x + 6 = 0, \qquad
2x^2 + 5x - 7 = 0, \qquad
x^2 + 4x - 5 = 0 .
\]
（提示：试取 $x = 1$ 与 $x = -1$，再利用两根之积。）
\end{exercise}

\begin{exercise}[$\star\star$]\label{exo:g11:quad:5}
求两个数，其和为 $14$，其积为 $45$。再求
周长为 $34$、面积为 $60$ 的矩形的长与宽。
\end{exercise}

\begin{exercise}[$\star\star$]\label{exo:g11:quad:6}
设 $f(x) = x^2 + (m-2)x + 1$，其中 $m$ 为实参数。当
$m$ 取何值时，\omterm{def:g10:algebra:equation}{方程} $f(x) = 0$ 恰有一个解？
无解？
\end{exercise}

\begin{exercise}[$\star\star$]\label{exo:g11:quad:7}
一个小球竖直向上抛出，经过 $t$ 秒后，它的高度为
$h(t) = -5t^2 + 20t + 1$（单位：米）。小球能达到的最大高度是多少，
在何时达到？小球在哪一时间\omterm{def:g10:numbers:interval}{区间}内高度至少为
$16$ 米？
\end{exercise}

\begin{exercise}[$\star\star$]\label{exo:g11:quad:8}
解\omterm{def:g10:algebra:equation}{方程} $x^4 - 13x^2 + 36 = 0$，令 $X = x^2$
（\emph{双二次\omterm{def:g10:algebra:equation}{方程}}）。
\end{exercise}

\begin{exercise}[$\star\star$]\label{exo:g11:quad:9}
对于怎样的 $x$ 值，点 $M(x, x^2)$ 位于\omterm{def:g11:quad:parabola}{抛物线} $y = x^2$
上，且严格处于直线 $y = x + 2$ 的下方？请结合草图说明。
\end{exercise}

\begin{exercise}[$\star\star$]\label{exo:g11:quad:10}
一个正数与它的倒数之和为 $\frac{13}{6}$。求这个
数。（化为一元二次\omterm{def:g10:algebra:equation}{方程}求解。）
\end{exercise}

\begin{exercise}[$\star\star\star$]\label{exo:g11:quad:11}
设 $\mathcal P$ 为\omterm{def:g11:quad:parabola}{抛物线} $y = x^2$，$d_m$ 为直线
$y = mx - 1$，其中 $m$ 为实参数。
\begin{enumerate}
  \item 当 $m$ 取何值时，$\mathcal P$ 与 $d_m$ 有
        两个交点？恰有一个交点？无交点？
  \item 当二者恰有一个公共点时，求这个切点。
        （在\cref{ch:g11:deriv}中，你将看出 $d_m$ 即为
        $\mathcal P$ 在该点处的切线。）
\end{enumerate}
\end{exercise}

\section{问题：抛物线的秘密点}

\begin{problem}[{周末问题——焦点与准线：为何
卫星天线、车头灯与悬索桥都画出同一种曲线，
而这样的曲线本质上只有一条}]
\label{pb:g11:quad:1}
每个阳台上的每一个卫星天线，都将其接收器对准
一个确定的内部点；每一个车头灯，都把灯泡藏在
同一个特殊位置。这个点——\emph{\omterm{pb:g11:quad:1}{焦点}}——是
\omterm{def:g11:quad:parabola}{抛物线}自希腊几何以来一直保守的秘密，而
距离公式便足以把它揭示出来。在此过程中，本题
将训练\omterm{def:g11:quad:discriminant}{判别式}，让抛射体越过墙壁，
悬挂起一座桥，并证明一个惊人的结论：在
相似意义下，\emph{世界上只有一条\omterm{def:g11:quad:parabola}{抛物线}}。

\medskip
\noindent\textbf{第一部分——\omterm{def:g11:quad:quadratic}{二次三项式}的熟练运用。}
\begin{enumerate}
  \item 依次解\omterm{def:g10:algebra:equation}{方程} $x^2 - 5x + 6 = 0$，$2x^2 + x - 6 = 0$，
        $x^2 + x + 1 = 0$（见\cref{thm:g11:quad:roots}）。
  \item 将 $2x^2 - 4x - 6$ 化为\omterm{def:g10:algebra:expand}{因式分解}形式与\omterm{thm:g11:quad:canonical}{顶点式}
        （见\cref{thm:g11:quad:canonical}）。连同
        \omterm{def:g10:algebra:expand}{展开}式，你现在便拥有同一\omterm{def:g10:functions:function}{函数}的三种表达形式：
        请一眼看出各形式分别显示何种特征
        （根、顶点、$y$ 轴截距）。
  \item 用符号规则解不等式 $-x^2 + 4x - 3 \geq 0$
        （见\cref{thm:g11:quad:sign}）。
  \item 韦达作为侦探（见\cref{prop:g11:quad:vieta}，
        延续\cref{pb:g10:algebra:1}中的巴比伦思路）：求两个
        数，其和为 $7$，积为 $12$，要通过写出它们的
        满足的\omterm{def:g10:algebra:equation}{方程}来求解。又知 $x^2 - 5x + c = 0$ 的一个根等于
        $2$，在不解任何\omterm{def:g10:algebra:equation}{方程}的情况下，
        求 $c$ 与另一根。
  \item 当 $m$ 取何值时，\omterm{def:g10:algebra:equation}{方程} $x^2 + mx + 9 = 0$ 有
        重根？无实根？有两个根？
\end{enumerate}

\medskip
\noindent\textbf{第二部分——\omterm{pb:g11:quad:1}{焦点}揭晓。}
研究\omterm{def:g11:quad:parabola}{抛物线} $y = x^2$、点
$F\left(0, \frac14\right)$ 与水平直线 $\delta$，其\omterm{def:g10:algebra:equation}{方程}为
$y = -\frac14$。
\begin{enumerate}[resume]
  \item 设 $P(x, x^2)$ 为\omterm{def:g11:quad:parabola}{抛物线}上一点，\omterm{def:g10:algebra:expand}{展开}
        $PF^2 = x^2 + \left(x^2 - \frac14\right)^2$，并证明它
        是完全平方 $\left(x^2 + \frac14\right)^2$。
        由此求出 $PF$。
  \item 计算点 $P$ 到直线 $\delta$ 的距离，
        并得出结论：\emph{\omterm{def:g11:quad:parabola}{抛物线}上每一点到 $F$ 的距离
        恰等于它到 $\delta$ 的距离}。（$F$ 是
        \emph{焦点}\index{焦点}，$\delta$ 是
        \emph{准线}\index{准线}。）
  \item 证明逆命题：若点 $P(x, y)$ 满足
        $PF = \operatorname{dist}(P, \delta)$，则必有
        $y = x^2$。（将两距离平方后化简。）
        \omterm{def:g11:quad:parabola}{抛物线}因此是一种\emph{轨迹}，恰如
        \cref{pb:g10:coordgeom:1}中的圆那样：
        由到一点与到一条直线等距确定，而非仅到一点等距。
  \item 对更一般的\omterm{def:g11:quad:parabola}{抛物线} $y = \frac{x^2}{4f}$（其中
        $f > 0$），仿照第 6 题证明其\omterm{pb:g11:quad:1}{焦点}为
        $(0, f)$，\omterm{pb:g11:quad:1}{准线}为 $y = -f$。应用：一个
        卫星天线宽 $60$ 厘米、深 $9$ 厘米——其
        边缘过点 $(30, 9)$。求 $f$：接收器必须悬挂在
        顶点上方多高处？
  \item 天线为何有效：与轴平行射入的光线
        经反射后都\emph{经过\omterm{pb:g11:quad:1}{焦点}}（位于\omterm{pb:g11:quad:1}{焦点}的灯泡
        发出平行光束，汽车前灯即是倒过来使用的
        天线）。解释第 7 题中的等距性为何
        使这一结论显得合理（可将光线想象为推进的波前），
        并说明将在\cref{ch:g11:deriv}中出现的哪一工具
        （见\cref{exo:g11:quad:11}）可将合理性
        变为证明。
\end{enumerate}

\medskip
\noindent\textbf{第三部分——\omterm{def:g11:quad:parabola}{抛物线}的应用。}
\begin{enumerate}[resume]
  \item 一个小球沿轨迹
        $y = x - \frac{x^2}{20}$（$x$、$y$ 以米为单位）运动。求
        小球落地的位置、最高点的\omterm{def:g10:coordgeom:system}{横坐标}，
        以及最大高度
        （见\cref{prop:g11:quad:variations}）。
  \item 同一轨迹上，在 $x = 4$ 处立有一堵墙。
        计算小球在该处的高度，判断它能否越过一堵
        $3$ 米高的墙？$4$ 米高的墙呢？
  \item 一个\omterm{def:g11:quad:parabola}{抛物线}形拱跨度为 $40$ 米，最大高度为
        $10$ 米。在以中心为原点的\omterm{def:g10:coordgeom:system}{坐标系}中，
        $y = 10 - \frac{x^2}{40}$。一辆高 $7$ 米的卡车能否
        在距中心 $10$ 米处通过？在 $15$ 米处呢？
  \item 某剧院票价为 $p$ 欧元时，其收入为
        $R(p) = p\,(120 - 2p)$（需求随价格上涨而下降）。
        求使收入最大的票价及其\omterm{def:g10:functions:extrema}{最大值}。
  \item 金门大桥的主缆在极高精度下呈\omterm{def:g11:quad:parabola}{抛物线}形
        （均匀的桥面荷载产生\omterm{def:g11:quad:parabola}{抛物线}；
        而自由悬挂的链条\emph{并非}\omterm{def:g11:quad:parabola}{抛物线}：其
        曲线即悬链线，需要用到十二年级的指数\omterm{def:g10:functions:function}{函数}）。
        已知跨度 $1\,280$ 米、下垂度 $143$ 米，
        以中心为原点，主缆曲线为 $y = 143\left(\frac{x}{640}\right)^2$。
        求主缆在最低点上方的高度，即
        $x = 320$ 米处比最低点高出多少？
\end{enumerate}

\medskip
\noindent\textbf{第四部分——直线、切线与最后的
惊喜。}
\begin{enumerate}[resume]
  \item 完成\cref{exo:g11:quad:11}：直线
        $y = mx - 1$ 过外点 $(0, -1)$，与
        $y = x^2$ 相交，其交点情况由\omterm{def:g11:quad:discriminant}{判别式}的符号决定。
        求从 $(0, -1)$ 出发的两条切线及其
        切点。
  \item 再考虑直线 $y = mx + 1$，它过\emph{内点}
        $(0, 1)$。计算其交点\omterm{def:g10:algebra:equation}{方程}的\omterm{def:g11:quad:discriminant}{判别式}，
        证明每条这样的直线都与\omterm{def:g11:quad:parabola}{抛物线}相交两次。
        由此得出结论：从内部看不存在切线
        ——正如圆的情形一样。
  \item \omterm{pb:g11:quad:1}{焦点}弦：将 $y = x^2$ 与过\omterm{pb:g11:quad:1}{焦点}的
        水平直线 $y = \frac14$ 相交，所得弦长是多少？
        检验一般规则（对天线值得记住）：
        对于 $y = \frac{x^2}{4f}$，过\omterm{pb:g11:quad:1}{焦点}
        且平行于\omterm{pb:g11:quad:1}{准线}的弦
        长为 $4f$。
  \item 最后的惊喜：施以缩放
        $(x, y) \mapsto (kx, ky)$ 作用于\omterm{def:g11:quad:parabola}{抛物线} $y = x^2$，
        证明其像恰为 $y = \frac{x^2}{k}$。由此得出结论：
        每条\omterm{def:g11:quad:parabola}{抛物线} $y = \frac{x^2}{4f}$ 都是放大了的单位
        \omterm{def:g11:quad:parabola}{抛物线}——不同于比例不同的圆和椭圆，
        \emph{所有\omterm{def:g11:quad:parabola}{抛物线}都相似}。
  \item 终章：用五句话概括\omterm{def:g11:quad:parabola}{抛物线}的全貌——
        一个有三种“装束”的\omterm{def:g10:algebra:equation}{方程}（第 2 题）、一个轨迹
        （第 7--8 题）、一个工程师的反射器（第 9、
        10 题），抛射体的轨迹与桥梁的缆索
        （第 11--15 题），以及一条可以缩放的单一曲线
        （第 19 题）。每项用一句话概括。
\end{enumerate}
\end{problem}
